% BEGIN TEMPLATE
\documentclass[11pt]{article}
\usepackage{graphicx}
\usepackage{hyperref} 
\usepackage{xcolor}
\usepackage{nameref}
\usepackage{listings}
\usepackage{float}
\usepackage[title]{appendix}
\usepackage[ruled]{algorithm2e}
\graphicspath{ {../../images/} }
\bibliographystyle{acm}
% CHANGE THESE
\newcommand{\courseListing}{CSCI 8360}
\newcommand{\courseName}{Machine Learning for Text}
\newcommand{\assignmentTitle}{Summary \#1}
\newcommand{\assignmentSubtitle}{What's Special About Machine Learning?}
\usepackage{geometry}
\geometry{margin=1in}

\hypersetup{
    colorlinks,
    linkcolor={red!50!black},
    citecolor={blue!50!black},
    urlcolor={blue!80!black}
}
\urlstyle{same}
\definecolor{codegreen}{rgb}{0,0.6,0}
\definecolor{codegray}{rgb}{0.5,0.5,0.5}
\definecolor{codepurple}{rgb}{0.58,0,0.82}
\lstdefinestyle{mystyle}{
    commentstyle=\color{codegreen},
    keywordstyle=\color{magenta},
    numberstyle=\tiny\color{codegray},
    stringstyle=\color{codepurple},
    basicstyle=\ttfamily\footnotesize,
    breakatwhitespace=false,         
    breaklines=true,                 
    captionpos=b,                    
    keepspaces=true,                 
    numbers=left,                    
    numbersep=5pt,                  
    showspaces=false,                
    showstringspaces=false,
    showtabs=false,                  
    tabsize=2
}

\lstset{style=mystyle}

\begin{document}
  \begin{center}
  \includegraphics[scale=0.15]{UNO-Logo-Color.png}
  \\[0.3in]
  \textbf{\courseListing{}}\\
  \courseName{}
  \\[0.75in]
  \textbf{\assignmentTitle{}}\\
  \assignmentSubtitle{}
  \\[0.75in]
  \textbf{Patrick Davlin}
  \\[0.75in]
  \textbf{Computer Science Department}\\
  \textbf{Peter Kiewit Institute}\\
  \textbf{University of Nebraska}
  \\[0.75in]
  \textbf{Fall 2021}
  \\[0.3in]
  \includegraphics[scale=0.075]{UNO-Icon-Color.png}
  \newpage
\end{center}
  \graphicspath{{./images/}}
  Analyzing text is an essential programmatic component of many online applications.
  Search engines, spam filters, recommendation algorithms, and social media site all rely on some amount of text mining in order to improve user outcomes for their products.
  This is hardly a surprising statement, because humans have primarily communicated via written word in some form or another since its invention by Sumerians.
  Underlying this obvious understanding is a more complicated discussion of how text is a difficult and nuanced domain for effective analysis.
  
  To illustrate this, it is easy to consider a counterexample covered in several other courses in the AI concentration coursework--image processing.
  Fundamentally, all images are objective--they have a defined size and are composed of individual pixels that display color in a set spectrum.
  The complexity of an image can be expressed in its file size--"how large, how detailed, how colorful is this image?"--and there are commonly-deployed strategies used to reduce the complexity of images for analysis by object recognition algorithms and neural networks.
  Recognition is, by extension, an analysis of this static set of information to determine whether groups of pixels constitute some specific shape or object in the image.
  
  Text, by comparison, occupies less "space"--its characters are encoded more simply--but contains complexity due to the combinations of characters that compose a document set (\textit{corpus}) and its complete set of words (\textit{lexicon}).
  Because sets of text can be small, mining has to occur across a large or massive corpora--many documents that might have different purposes, writers, and styles.
  This also requires accounting for the variety of lengths that documents in a corpus may have, to avoid emphasizing individual words too much.
  
  Depending on how the text is analyzed--as a pure "bag of words" or as groups of phrases in context with a language model--applications can introduce substantial dimensionality into their algorithms, possibly in the range of $10^5$ or $10^6$ dimensions, according to the textbook.
  The meaning of an individual word can differ depending on context, and the most important words might only appear once or twice in a corpus, which makes text a very \textit{sparse} medium.
  Broken down into dimensions, many words have "zero value"--they do not contribute to the overall meaning of the text--but some dimensions carry positive value.  
  
  Language models "create a probabilistic representation of the text" by trying to capture the likelihood that some word might be used in a particular sequence.
  Other approaches create topic modeling or clustering based on the "value" of a word.
  Since most dimensions in a model carry a zero value, the positive value words are highly valuable sources of information for models.
 
  Not emphasized too much in the textbook but still inherent to the medium of text is the prevalence of individual writers' usage--or mis-usage--of language, including misuse of words, misspellings, slang terms, or strange linguistic tricks like portmanteaus or hyperbole.
  Different people, in varying regions, from diverse backgrounds, with a range of backbrounds, and across all generations, can use the same language (or even the same lexicon) in a tremendously varied manner.
  Also, for technical documents, extraneous information also needs to be considered for removal, like tag elements in HTML, or categorization, such as the name of a user sending a message in a forum.
  
  Text analysis is difficult because the usage of text in any language is unique.
  Some applications, like the famous Google search engine and many recommendation algorithms, have very advanced text mining and analysis approaches.
  One interesting opportunity that a student in this class could hope for is to better understand the methods that some successful everyday tech products use to work so well.
  In a more academic (and wishful) sense, having a better understanding of the analysis of language might also provide better insight into language itself--and provide a student with a better opportunity to understand the world they live in.
\end{document}